% !TEX root = ../PhD Thesis.tex
\chapter*{Eulogy}
\addcontentsline{toc}{chapter}{Eulogy}

 I leave the most important acknowledgment to my grandparents and parents. The lack of an Oxford comma is intentional. They were, effectively, both. Born into a couple in the midst of divorce, I was raised by my maternal grandparents through the early years of my life, and I was deeply shaped by them. I carry both their surnames, and far more of them in me than I can name. They gave me the greatest support of my life, filled with happiness and laughter, and inevitably the greatest heartache. It feels as if my heart has been frozen ever since. Every time I allow myself to feel, that ice shatters all over again, and the shards stab with each heartbeat. Only by not allowing myself to feel, by focusing on anything else, running from my own mind that so insistently gravitates toward traumatic thoughts, can it freeze again, granting me some respite. The constant frostbite burning is after all familiar, and less intense than the stabbing. They both passed away to the same illness, the one I have dedicated my career to understanding. It gives me purpose, and a constant anger that always pushes me forward, no matter how scalding the sun on this path is. But that fuel is not without its price. Every failure becomes ever more personal. Every day of delay, more deaths that perhaps could have been prevented. Perhaps if I had worked harder, someone could have been spared their fate, or I could have been spared my own, the powerlessness of watching loved ones fade away. Every small step down a Kaplan-Meier curve is someone who likely suffered just as much. Every study, everything I learn about this disease, inevitably ties back to those experiences. I do not know if I will ever achieve a meaningful contribution. Even this work, into which I poured so much of myself, will likely leave no mark. I found errors, and at most identified a biomarker whose impact is, all things considered, modest. Nevertheless, if I ever do achieve something, it will undoubtedly be because of them. And so I give my thanks.
 \\
 \\
 To my grandfather, my father, my best friend for so long. Someone who stayed positive no matter how bleak things looked, someone, in that way, the opposite of me. Always the central pillar of our family. We played games together, and the things you loved rubbed off on all of us, even in the smallest ways. When I first moved to another city to study, it became tradition to call each other whenever our team had a good result. After so many bleak years, we rose so high. If only you could have seen. We were champions twice over, made remarkable runs in the Champions League, reaching the final knockout stages. In those last years we beat Tottenham (before they fell apart), beat Arsenal, hammered the all-powerful Manchester City five to one, and won against the reigning European champions. In every one of those games, the joy of winning was always followed by the impulse to call you, only to be followed by the heartache that that is no longer possible. All of us, children and grandchildren alike, still call one another, but we all feel it: that void, that sense of something missing. It is just a sport, one that is given far too much importance in the face of so many pressing issues, especially in a country like ours. But because of all the memories we made over it, it will always be more than a game to me.
 \\
 \\
 How is it that even the smallest things can be so bound to a person, carrying such weight? I remember once mentioning offhandedly that I needed toothpaste when I got back to Lisbon. A few moments later you gave me one. When I arrived, I discovered I already had a full supply. And yet, even after that supply finished, I could not use that one for more than a year, simply because it had come from you. I still keep many memories in those forms, pens you used, your necklace, other small objects. Here, I used to watch almost every match with a friend. After you were gone, even as we were on the verge of becoming champions again, after just a short break, I could not watch any of it. They were too bound to you. When we were champions once more, I was given a shirt, and I had our surname put on the back, as a tribute to those memories. Years have passed now, and the pain has not eased. Every achievement, no matter how much joy it brings, will always bring pain alongside it. I cannot enter a hospital without thinking of those last terrible moments , the methodical sound of the ventilators, a person so bright reduced to a shadow, and I do not think I ever will.
 \\
 \\
 To my grandmother, my mother, the person I admired most for so long. Where his illness struck suddenly, yours was slow. And likewise, so much of what I am came slowly from you. I changed my name to carry yours. I will never understand the logic of how surnames are inherited. The name I chose technically came from your grandfather, but I chose it because of you. You always impressed upon me the importance of studying, of working hard, of learning how to think, of how much it mattered for the future. You were always curious, always eager to discover new things. It is more than fitting that my publishing name is yours , the name of the person who most shaped my own curiosity and desire to learn. We both always loved history, and we shared things we had found: obscure facts that most people would never find interesting. I remember reading that our own King Peter the Cruel, heralded as a great romantic, had another mistress shortly after Inês, and was, all things considered, a bit of a bastard. I wanted to share that with you, to break apart that preconceived notion. Or the stories of great women who achieved so much against the weight of their time: Tamar of Georgia, Urraca of Castile, Matilda of Canossa, Eleanor of Aquitaine, so many things I wanted to share, and then you were no longer there. That passion still holds. I still fall into deep holes of research on obscure historical subjects, but always with this sadness underneath, the knowledge that the person I shared it with, the person who gave it to me, is gone.
 Perhaps I carry more regrets about you than I can properly account for. Your illness began earlier, and a few years in I enrolled in oncobiology, driven by the desire to help you in whatever way I could. During my training, it became clear that the treatment you were receiving was not the right one. I contacted others who agreed, and we managed to change it. I believe you lived a few more years because of that, but I will never fully come to terms with whether it was the right thing to do. Because of me, you lived long enough to watch your husband die an atrocious death. That is not what I would want for myself, and I do not think I will ever resolve it. The desire to do great things so often prevents us from simply doing good.
 \\
 \\
 I was never squeamish. During my degree in biomedical sciences we observed surgeries, and the sight of blood, though uncomfortable, was something I could manage, which is ironic, given that I have always fainted when giving blood. But now I have seen enough blood for a lifetime. A bathtub full of it, requiring many transfusions to replenish. You were always proud of my achievements, but you were already too ill to see the master's defence.
 \\
 \\
 Both of you were always supportive of what I did. I wish with everything I have that you could see me defend this thesis. I will always carry the regret that I was not fast enough to give you that. You were so much of me, and I am so much less than myself ever since your passing.